\begin{figure*}[t]
\centering
\includegraphics[width=\textwidth]{Figure_B1_fixed_task0_gradient_radius50_reference_style.pdf}
\caption{Fixed Task0 positive-gradient displacement diagnostic using formal ImageNet-A seeds 0, 1, and 2 (ViT-B/16, rank-10 Q/V LoRA). Each method/seed uses its own Task0 checkpoint and one gradient of the mean cross-entropy over all 1,259 Task0 training examples. Normalize this gradient once and evaluate theta(r) = theta0 + r*u, keeping the direction fixed across all radii. Perturbations and gradients use the 24 effective Q/V matrices; the old-task 20-class head is fixed. Left: signed training CE change relative to the corresponding anchor: for method m and seed s, D(m,s,r) = L(m,s,theta(m,s,r)) - L(m,s,theta(m,s,0)). The overbar denotes the arithmetic mean of these changes across three seeds within each method, with no positive clipping. Middle: full Task0 training accuracy (\%). Right: raw empirical Fisher trace recomputed at each position on the same fixed 256 training examples, as the mean squared norm of per-example true-label CE gradients. There is no extra Fisher normalization or multiplication by lambda. Lines show seed means and shading shows one sample standard deviation (ddof=1); the independent replication unit is the training seed, not the direction or radius. Directions differ between methods. Connecting lines are visual guides between measured points. This is a fixed-direction diagnostic, not an actual future-task learning trajectory or a guarantee of monotonic loss at large radius. The figure displays all 93 measured radii from 0 to 50; the separate 0–2 companion resolves the initial response. }
\end{figure*}
